\chapter{Background and Formal Query Model}
\label{chap:background}

\section{Related Work}

\subsection{SUQL}
SUQL~\cite{suql} extends SQL with a free-text predicate function,
written \answerfn{}$(\mathit{text}, \mathit{question})$, that is backed
by an LLM and returns a boolean (optionally with a short justification).
Its central systems argument, which this project adopts wholesale, is
that \emph{predicate pushdown} -- executing every structured predicate
\emph{before} any semantic predicate -- is the dominant lever for
controlling LLM invocation cost, because it shrinks the candidate set
the expensive operator ever has to see. SUQL treats \answerfn{} as an
opaque, per-row, non-batchable user-defined function (UDF); it does not,
in its base form, address the cost of the semantic operator itself once
the candidate set is fixed. That is exactly the gap this project's
cascading contribution targets.

\subsection{ELEET}
ELEET~\cite{eleet} generalizes the single-source assumption of SUQL: a
logical scan is modeled as a \emph{multi-modal operator} that can pull
evidence for the same entity from more than one physical source --
structured columns in one table, free text in another -- and fuse it
before a selection is applied. This is directly relevant to our setting,
where the IMDb structured metadata and the review text are joined from
two logically separate sources rather than pre-materialized into one
table. ELEET also formalizes an \emph{operator cascade}: a sequence of
selection operators, each with an associated cost and selectivity, that
can be reordered to minimize expected total cost, and each parametrized
by a \emph{log-odds} score rather than a hard boolean, which is the
representation this project's cascade calibration is built on
(Section~\ref{sec:log-odds}).

\subsection{Stretto}
Stretto~\cite{stretto} is the most recent piece of prior work this
project draws on, and contributes two largely independent ideas:
\begin{enumerate}
  \item \textbf{Bayesian threshold calibration.} Rather than hand-tuning
        the log-odds cutoff at which a semantic operator's decision is
        trusted without further verification, Stretto derives a
        \emph{credible interval} on the operator's precision and recall
        at any candidate threshold from a Beta posterior fit to a
        calibration set, and picks the threshold(s) that certify a
        target quality floor at a target confidence level. This
        procedure is a pure post-processing step over whatever log-odds
        score the cheap stage already produces; it requires no change
        to how that score is computed, and is therefore
        \textbf{backend-agnostic}.
  \item \textbf{Cache-aware execution ladder (Expected Attention Press).}
        Stretto's second contribution reorders and prunes cascade
        stages at a much lower level: it reuses attention key/value
        (KV) cache state \emph{across} cascade stages, so that escalating
        an ambiguous item from a cheap to an expensive stage does not
        require recomputing attention over the shared context from
        scratch. This requires direct access to the serving engine's
        attention state (\eg via \texttt{vllm}'s internal APIs), and is
        \textbf{not portable} to an Ollama-served backend without
        replacing the serving layer.
\end{enumerate}
This project builds on contribution (1) only; contribution (2) is
deferred to future work (Section~\ref{sec:future-kv-cache}) pending a
migration away from Ollama. Neither Stretto nor the systems it builds on
state a formal condition for \emph{when} a calibrated cascade is
guaranteed to reduce cost relative to an all-expensive baseline;
Section~\ref{sec:theory-cost} closes this gap for the two-level cascade
studied here, with a finite crossover point and a concentration
guarantee, instantiated against measured system constants.

\subsection{Positioning}
Table~\ref{tab:related-positioning} summarizes how this project's
contributions relate to the three pieces of prior work above.

\begin{table}[h]
\centering\small
\begin{tabular}{p{3.1cm}p{9.5cm}}
\toprule
Source & What we adopt / extend \\
\midrule
SUQL & Predicate-pushdown discipline; \answerfn{} as the expensive,
       per-row semantic operator inside the cascade's second stage. \\
ELEET & Multi-modal operator abstraction for the two-source join;
        log-odds parametrization of a semantic filter; operator-cascade
        cost/selectivity reordering rule (Section~\ref{sec:operator-ordering}). \\
Stretto & Bayesian Beta-posterior calibration of the cascade's
          escalation band (Section~\ref{sec:beta-posterior}); we do not
          (yet) adopt the KV-cache execution ladder. \\
\bottomrule
\end{tabular}
\caption{How this project's cascaded design relates to prior work.}
\label{tab:related-positioning}
\end{table}

\section{Formal Query Model}
\label{sec:formal-model}

We fix the following notation, used throughout the rest of the report.

\begin{itemize}
  \item $R$: the structured relation, schema
        $(\mathit{id}, \mathit{year}, \mathit{genre}, \mathit{rating},
        \mathit{cast}, \dots)$.
  \item $D$: the unstructured relation, schema
        $(\mathit{id}, \mathit{review})$.
  \item $\sigma_S$: a structured predicate over $R$'s columns,
        compilable to a SQL \texttt{WHERE} clause.
  \item $\sigma_U$: a semantic predicate over $D$'s free text, expressed
        in natural language and evaluated via \answerfn{}$(\mathit{review}, \sigma_U)$.
  \item $R_{\sigma_S} = \{\, r \in R : \sigma_S(r) \,\}$: the structurally
        filtered relation.
  \item $s = |R_{\sigma_S}| / |R|$: the \emph{selectivity} of the
        structured predicate.
  \item $\mathrm{gold} \subseteq R \bowtie D$: a hand-verified gold
        answer set for a given query $Q$, used to compute precision,
        recall, and F1.
\end{itemize}

The target output of any execution strategy is
\[
  \widehat{\mathrm{ans}} \;=\; \bigl\{\, r \in R_{\sigma_S} \bowtie D \;:\; \sigma_U(r) \;\text{holds} \,\bigr\},
\]
and every strategy studied in this report is judged on how closely
$\widehat{\mathrm{ans}}$ tracks $\mathrm{gold}$ (precision/recall/F1), and
on how many LLM calls and how much wall-clock time it takes to produce
$\widehat{\mathrm{ans}}$.

\subsection{Log-odds representation of a semantic decision}
\label{sec:log-odds}
Any operator (cheap or expensive) that evaluates $\sigma_U$ against a
unit of work $u$ (a row or a batch of rows) is modeled as returning a
probability $p = f(u, \sigma_U) \in (0,1)$ that $u$ satisfies $\sigma_U$,
rather than a bare boolean. We work in \emph{log-odds} space,
\[
  \ell = \operatorname{logit}(p) = \log \frac{p}{1-p} \in \mathbb{R},
\]
for two reasons that matter operationally, not just aesthetically:
\begin{itemize}
  \item \textbf{Additivity.} Under (approximately) independent evidence
        combination, log-odds from multiple signals add, which makes it
        natural to combine a cheap-stage score with, \eg, a
        self-consistency vote or an embedding-similarity prior, without
        renormalizing probabilities by hand at every step.
  \item \textbf{Symmetric decision point.} $\ell = 0$ corresponds to
        $p = 0.5$, the natural undecided point, and $|\ell|$ is a
        monotone measure of confidence in either direction, which is
        exactly the quantity the cascade thresholds on
        (Chapter~\ref{chap:cascading}).
\end{itemize}
